\section{Balance sheet analysis}

\begin{qfmethod}[Intangibles, goodwill and the quality of reported equity]\label{meth:intangibles-goodwill-analysis}
\qffield{Statement} An identifiable intangible is separable or arises from contractual/legal rights. Purchased intangibles enter at cost; identifiable intangibles acquired in a business combination are separately measured at acquisition-date fair value, even if absent from the seller's balance sheet. A finite-lived asset is amortized from availability for use and tested for impairment when indicated; an indefinite-lived asset is not amortized and is tested annually. ``Indefinite'' means no foreseeable limit, not immortal. Internally generated goodwill is not recognized \citep[IAS~38 paras.~12, 21, 33, 48, 97, 107 and~118; IAS~36 para.~10]{adopted-accounting-standards-2023}.
\qffield{Practice} In an isolated acquisition of all shares, consideration 150 buys identifiable assets at fair value 180 and liabilities 60. Goodwill is $150-(180-60)=30$, assuming no deferred-tax complication, prior holding or noncontrolling interest. Reclassifying a valuable customer contract into goodwill improperly would avoid amortization and inflate subsequent earnings. Under IFRS, goodwill is allocated to cash-generating units and impairment is allocated first to goodwill; it is never reversed. US public-company goodwill uses reporting units and a quantitative loss equal to carrying amount above fair value, capped at goodwill, rather than the obsolete implied-goodwill second step \citep[IFRS~3 paras.~10, 18, 32; IAS~36 paras.~104 and~124]{adopted-accounting-standards-2023}; \citep[Summary pp.~1--3]{us-goodwill-2017}.
\qffield{Interpretation} Inspect goodwill by business, acquisition assumptions, useful lives, amortization, impairment drivers and sensitivity disclosures. Excluding goodwill from tangible equity may help assess collateral or acquisition dependence; it does not prove goodwill has zero economic value. IFRS revaluation of an intangible requires an active market, rarely available for unique brands; US GAAP generally does not permit upward revaluation. Private-company goodwill alternatives are outside this public-company example.
\end{qfmethod}

\begin{qfmethod}[Financial instruments, debt measurement and liability classification]\label{meth:instruments-liabilities-reporting}
\qffield{Statement} IFRS 9 debt assets combine a business-model test with contractual cash flows consisting solely of principal and interest: hold-to-collect supports amortized cost, collect-and-sell supports FVOCI, and otherwise FVTPL applies, subject to designation rules. Credit-loss allowances remain relevant. Equity investments default to FVTPL; an eligible non-trading equity investment can elect irrevocable FVOCI, with no recycling of its value changes to profit on sale. US debt securities distinguish held-to-maturity, available-for-sale and trading, while equity securities within ASU 2016-01 generally carry fair-value changes through income, with a measurement alternative for eligible investments without readily determinable fair values. Neither framework labels every security ``available for sale'' \citep[IFRS~9 paras.~4.1.1--4.1.4 and~5.7.5, B5.7.1--B5.7.1A]{adopted-accounting-standards-2023}; \citep[Summary pp.~1--3]{us-instruments-2016}.
\qffield{Practice} A three-year bond paying 5 annually and 100 principal costs 97.3270 at an effective rate of 6 percent. First-year effective interest is 5.8396; cash interest is 5, so amortized cost rises to 98.1666. If its year-end fair value is 102, an eligible FVOCI debt asset has an additional 3.8334 valuation gain in OCI (ignoring credit/FX effects), while FVTPL includes that gain in income. The same effective-interest logic amortizes a borrower's discount; reporting net debt after issuance costs does not reduce contractual principal due. US ASU 2015-03 presents issuance costs of recognized debt as a deduction from the debt liability, not a separate asset \citep[Summary pp.~1--2]{us-debt-costs-2015}.
\qffield{Conditions} Inspect maturities, covenants, collateral, defaults, fair values, measurement categories and liquidity/credit/market-risk disclosures. Under IAS 1's amendments effective from 2024, covenants required only after the reporting date do not change classification at that date, but require disclosure of the risk of early repayment. A covenant breach at the reporting date can change classification; evaluate the actual right to defer settlement and waiver timing. IFRS 9/7 amendments effective from 2026 clarify contingent cash-flow features and electronic-settlement derecognition and add disclosures; plain fixed-coupon examples here do not model those optional features \citep[effective-date and covenant paragraphs]{liability-covenants-2024}; \citep[classification, settlement and effective-date paragraphs]{financial-instrument-amendments-2026}.
\end{qfmethod}

\begin{qfmethod}[Common-size balance sheets and the liquidity-solvency distinction]\label{meth:common-size-balance-sheet-analysis}
\qffield{Statement} Divide balance-sheet items by total assets to compare financing and asset composition. The shares on each side sum to 100 percent, but netting, consolidation and measurement choices affect those shares. Liquidity concerns near-term payment capacity; solvency concerns the longer-run burden of obligations. Current and quick ratios use liabilities due within the operating cycle or the next year, while debt/equity must specify whether ``debt'' means interest-bearing debt or all liabilities \citep[chap.~4, solutions~18 and~25--34, PDF pp.~182--183]{statement-analysis-workbook-2020}.
\qffield{Practice} Meridian's 2026 assets total 867.5. Cash, receivables, inventory and PP\&E are 10.09, 17.29, 23.05 and 49.57 percent. Payables, accruals, debt and equity are 13.83, 5.19, 25.36 and 55.62 percent. Current assets are 437.5; current liabilities are $120+45+30=195$, including debt due next year. Current ratio is 2.244, quick ratio $(87.5+150)/195=1.218$, and cash ratio $87.5/195=0.449$. Interest-bearing debt/equity is $220/482.5=0.456$; total liabilities/equity is $385/482.5=0.798$.
\qffield{Interpretation} Cash's asset share fell from 12.5 percent and receivables grew faster than sales. Leverage nevertheless declined: debt/equity was $250/410=0.610$. The company has improved solvency measures but weaker immediate cash cover, a useful counterexample to treating every ratio as the same signal. Inventory obsolescence or disputed receivables could make a current ratio above two less reassuring than it appears. Check maturity timing and usable cash, not just a year-end snapshot.
\end{qfmethod}
